\documentclass[11pt]{article}

\usepackage[margin=1.2in]{geometry}
\usepackage{amsmath,amssymb}
\usepackage{booktabs}
\usepackage{url}
\usepackage[hidelinks]{hyperref}
% Document properties a reader sees when the file is downloaded. The journal
% build generates its own preamble and carries none of this.
\hypersetup{
  pdftitle={Auditing the ARC-AGI-3 Public Benchmark: Human Baselines, Scoring Rule, and Environments},
  pdfauthor={Ibrahim Mian and Shayaan Siddique},
  pdfsubject={An audit of the public ARC-AGI-3 reference artefacts},
  pdfkeywords={ARC-AGI-3, benchmark audit, formal verification, differential testing}
}

\newcommand{\code}[1]{\texttt{#1}}

\title{Auditing the ARC-AGI-3 Public Benchmark:\\
Human Baselines, Scoring Rule, and Environments}
\author{Ibrahim Mian \qquad Shayaan Siddique\\[2pt]
\normalsize Millennium Research\\[2pt]
\normalsize \texttt{\{ibby,shayaan\}@millenniumresearch.ai}}
\date{}

\begin{document}
\maketitle

\begin{abstract}
ARC-AGI-3 scores an agent by its action count against per-level human
baselines, so its results are only as sound as those baselines, the conditions
under which they were measured, and the code that compares agents with them.
We audit all three for the 25 public environments: we read the released human
replays against the shipped code, and check the scoring rule, the benchmarking
harness and two environment levels with machine-verified models run
differentially against the implementation.

The 340 released recordings reproduce 165 of the 183 published baselines
exactly as the upper median of per-play action counts with resets charged, and
on the 72 levels that tell the charged and uncharged readings apart the
published value is the charged one on 65: humans paid for their resets as
agents do. Of the other 18 baselines, two carry an earlier recorder's tally, one
action low, and 16, in three environments, are not reproduced by the release.
Most baselines rest on few plays: under resampling, the median baseline's 95\%
interval is 84.6\% as wide as the baseline itself. All 340 recordings reproduce state and level at every one of their
$180{,}496$ steps through the shipped environments, 19 of them only under an
engine switch the default does not set: on the engine the humans played, a
reset never leaves the level, while on the shipped default a reset at the start
of a level, or straight after another, returns the player to level 1. Every
published baseline is at or above its level's optimum.

In the harness, the reset it forces after a game over consumes the agent's
per-level budget, so the effective budget is $\lceil 5b \rceil$ minus the number
of deaths, and a paired counterfactual shows this can end a level on a real
public environment. Its reach is small: applied to the 1,614 level completions
in the human replays, the budget cuts off 29, none of them because of those
resets, and no finding in this paper is shown to change a published score. We also report an under-specified aggregation rule, of which the
toolkit implements one reading, and an equation in the technical report that
permits a per-level score of $132.25\%$ where the report's own stated range,
the documentation and the code cap it at $115\%$. Two environment models agree with the shipped implementation on all
$190{,}560$ transitions of their levels. The code, models, preregistrations and
derived data are public at pinned commits.
\end{abstract}

\section{Introduction}

ARC-AGI-3 \cite{arcagi3}, a successor to the Abstraction and Reasoning
Corpus \cite{chollet}, is an interactive-reasoning benchmark: an agent is
placed in a turn-based environment whose rules are stated nowhere, and must
discover them. Scores are efficiency ratios. For each level of each environment,
the agent's action count is compared against a human baseline, the ratio is
squared and capped, the level scores are weighted by level index, and the
environment scores are averaged. The benchmark is unusual among evaluation
suites in how much of it is public: the toolkit, the harness, the example
agents, 25 of the 135 environments, and a technical report describing the
scoring rule in detail.

That openness makes an audit possible, and an audit worthwhile. A score under
this rule depends on three artefacts being what they claim: the environments,
the human baselines that agents are measured against, and the scoring code. This paper
checks all three, as far as each admits checking.

Much of what decides a score here is data rather than code. Every level score
is a ratio against a number measured from people, and the comparison between
agents and humans rests on the two having been measured the same way. The
human baselines, the recordings they were computed from, and the conditions
under which those recordings were made are therefore as much a part of the
benchmark as its environments. Section~\ref{sec:baselines} audits them as data:
it reproduces the published baselines from the released recordings, identifies
the rule that produced them, and replays every recording through the shipped
code to recover the conditions under which the humans played.

Our position is not adversarial. Benchmarks are infrastructure, and the useful
service an outside reader can perform is to say precisely which parts have been
checked, by what method, and what survived. Where we found nothing wrong we say
so at the same length as where we found something.

\paragraph{Contributions.} None of the findings below is shown to change a
published score; each is reported with the reach we could measure.
\begin{enumerate}
\item An audit of the published human baselines as data
      (Section~\ref{sec:baselines}). The 340 released human replays reproduce
      165 of the 183 baselines exactly under one stated rule, settle that
      human resets were charged, show every baseline at or above its level's
      optimum, and leave 16 levels in three environments that the release does
      not reproduce, seven of whose published values occur in no released play.
      Resampling the plays shows how noisy the baselines are: the median
      baseline's 95\% interval is 84.6\% as wide as the baseline itself.
\item The conditions under which the baselines were collected
      (Section~\ref{sec:humanengine}): the humans played under reset semantics
      equivalent to the engine's \code{ONLY\_RESET\_LEVELS} switch, on which a
      reset never leaves the level, while the shipped default turns a reset at
      a zeroed action counter into a full reset to level 1. All 340 recordings
      reproduce state and level through the shipped code under that switch,
      and 19 only under it, exactly the plays that contain such a reset beyond
      level 1. A harness probe shows the same 59 actions scoring $10.71$ with
      two levels under the humans' setting and $3.57$ with one under the
      default.
\item A finding in the benchmarking harness (Section~\ref{sec:forced}): it
      spends the agent's per-level budget on the resets it issues after a game
      over, so the effective budget is $\lceil 5b \rceil$ minus the number of
      deaths, which ends a level the same policy otherwise completes on a real
      public environment at its real budget. The exposure is measured across
      all 183 public levels and the reach on human play across all 1,614 human
      level completions.
\item A machine-verified model of the published scoring rule, run
      differentially against the shipped scorer on preregistered planted traces
      (Section~\ref{sec:scoring}). On every planted trace with an oracle value
      but one, the scorer agrees with our model of the rule as the report's
      prose states it, and the one disagreement is an unreached edge; five
      documentation problems are reported: the equation's cap, the budget's
      location, the unstated reset accounting, a stale code comment, and an
      omitted schema field.
\item Machine-verified models of two levels of a public environment, agreeing
      with the shipped implementation on every transition of their complete
      reachable state graphs (Section~\ref{sec:env}), with exhaustive state
      counts for those levels under a stated identity function, checked against
      the report's own stated win probability; and reset and level-accounting
      probes covering all 25 environments (Sections~\ref{sec:sweep}
      and~\ref{sec:play}).
\end{enumerate}

\section{Scope, and what we did not do}
\label{sec:scope}

In scope: the 25 public demonstration environments, the toolkit
\code{arcprize/ARC-AGI} at commit \code{f12822c}, the harness
\code{arcprize/arc-agi-3-benchmarking} at \code{1aa78da}, the example agents
\code{arcprize/ARC-AGI-3-Agents} at \code{4743e7d} (all MIT-licensed), the
technical report, and the public documentation.

Out of scope, deliberately and throughout: the 55 semi-private and 55 private
environments. We made no attempt to recover, probe, scrape or infer them. All
environment execution was local and offline. No component of this work is an
agent whose purpose is to score on the benchmark. The random player of
Section~\ref{sec:play} probed reset semantics and won no game; the blind policy
of Section~\ref{sec:forced} measured a cost and completed a level in 16 of its
250 runs without attempting to. Nothing was submitted to any leaderboard.

We audit the \emph{shipped local} artefacts at the pinned commits. The
server-side scorer that produces the official leaderboard is not observable to
us, and we make no claim about it.

\section{Method}
\label{sec:method}

Four commitments shape everything below.

\paragraph{Reject-only.} No automated component of this work emits a verdict of
correctness. Each either reports a disagreement or reports that it found none,
and a human reads every disagreement and classifies it before it is called
anything. ``No disagreement on $N$ fixtures'' is a statement about $N$ fixtures.

\paragraph{Preregistration.} The bar, the predictions, and the reading of each
possible outcome were written down and committed before each measurement.
The preregistration was committed (\code{8e0f70b}) before the first measurement
(\code{08b4f2e}), and each later phase's plan precedes its runs in the same
history; commit order in a history we control is evidence of sequence, not
proof, and we state it as such. Deviations are dated in a decisions file and disclosed. One preregistered
prediction was wrong and is reported as such in Section~\ref{sec:scoring}.

\paragraph{Machine-verified models, run differentially.} For each rule under
audit we write a model in Dafny \cite{dafny} whose properties are proved rather
than tested, compile it to JavaScript, and compare it against the shipped
implementation on traces recorded from that implementation. A model carries no
\code{assume}, \code{\{:axiom\}}, \code{\{:verify false\}} or \code{\{:extern\}},
and a proof obligation that times out is treated as unknown, never as
discharged. In the models of Section~\ref{sec:env} the level constants are
extracted from the shipped level data by a script, so a hand edit to them
fails the test suite; the transition rules are hand-written and are checked
only by the differential comparison.

\paragraph{Claims discipline.} Every number in this paper is produced by a
script into an artefact and registered in a claims file, and a checker verifies
that each value appears both in the text and on a matching line of the log that
produced it. The repository ships the scripts, the artefacts and the checker.

\section{The human baselines}
\label{sec:baselines}

Throughout, a \emph{cell} is one level of one environment, 183 in all, and the
\emph{upper median} of $n$ sorted counts is the one at position
$\lfloor n/2 \rfloor$ counting from zero, the statistic the report names for its
baselines.

\subsection{A check that needs no human data}

Every ARC-AGI-3 score is a ratio against \code{baseline\_actions[l]}. A human
cannot complete a level in fewer actions than the level's optimum, so a
published baseline must be at least that optimum. Breadth-first search supplies
the optimum, and supplies a sound bound on it even when it does not finish:
having expanded every state at depth $d$ without reaching a win, it has proved
no solution exists in $d+1$ actions or fewer.

Two things were settled before measuring, either of which would have made a
contradiction spurious. The scorer pairs \code{baseline\_actions[level\_idx]}
with that level's own action count, computed as a difference between cumulative
counts, so the published values are per level. And the published values are not
monotone --- \code{ar25} is $[32, 50, 75, 37, 89, 159, 233, 73]$ --- so they
cannot be cumulative.

\subsection{Result}

Of 18 levels attempted across the six enumerable environments (Section~\ref{sec:enumerable}), six resolved and
all six are consistent; none is impossible; twelve were stopped by a time or
memory cap after completing depths between 8 and 16 against baselines between 26
and 183, and establish nothing either way.

\begin{table}[t]
\centering
\begin{tabular}{llrr}
\toprule
Environment & Level & Minimum actions & Published baseline \\
\midrule
\code{ls20} & 1 & 13 & 22 \\
\code{ls20} & 2 & 45 & 123 \\
\code{ls20} & 3 & 39 & 73 \\
\code{tu93} & 1 & 18 & 19 \\
\code{tu93} & 2 & 10 & 16 \\
\code{tu93} & 3 & 19 & 34 \\
\bottomrule
\end{tabular}
\caption{Exact optima against published human baselines, for the levels where
the search completed. Every optimum is replayed through the shipped game and
shown to complete its level.}
\label{tab:baselines}
\end{table}

Table~\ref{tab:baselines} gives the resolved levels. The tightest margin is
\code{tu93} level 1, where the upper-median human used one action
more than the optimum on an environment whose rules are stated nowhere.

We then attempted to resolve the remaining levels by changing the search, and
could not. Breadth-first holds a whole layer of game objects, so memory was the
obvious suspect; we added depth-limited depth-first search with shallowest-depth
memoisation, which gives the same completeness guarantee to a stated depth while
holding only the objects along one path, and an iterative-deepening form that
banks a bound after each completed depth. Memory fell by roughly a factor of
seven on \code{g50t} level 1, where we measured it. It resolved nothing: in the same time budget
breadth-first reached a greater depth on both levels compared, because deepening
re-explores, and on levels breadth-first can finish it is strictly better. The
binding constraint is not memory but the size of the state space at the depths
these baselines occupy, and on those twelve levels neither exhaustive search we
tried reaches depth 26 within the caps of Appendix~\ref{sec:caps}, let alone
183. The human replays settle them from the
other side (Section~\ref{sec:replays}): a human's count is an upper bound on the
optimum, and every one of the twelve has a human count at or below its
baseline.

These optima are proved against the shipped local environment at the pinned
version. The replays show that the humans played that code, under one engine
switch the default does not set (Section~\ref{sec:humanengine}).

\subsection{The replays, and an error of ours}
\label{sec:replays}

An April 2026 post \cite{humandataset} states that the ARC Prize Foundation open-sourced
the Public Demo dataset, ``which includes 342 human step-by-step replays'' across
the 25 public environments. An earlier version of this work, published in its
repository, said we could not locate them. That was wrong, and the error was ours. Our scripted check of the
Foundation's public channels found no replay dataset among the \code{arcprize}
GitHub repositories or Hugging Face datasets, and the one link the announcement
gives, a link shortener, refused our request with HTTP 403 whether or not it
carried a browser user-agent; we reported the refusal as the link not resolving.
Two automated blind reviews of that version pointed out that a 403 from a
shortener is what a scripted request gets. Opened in a real browser on 5 September 2026, the link
resolves to a Google Drive folder holding the dataset: a folder with one
subfolder per public environment, an archive of the same, and a ratings file,
dated 13 to 15 April 2026. The observation is recorded as such in the
repository, with its provenance stated.

We record the error rather than deleting the trace of it, because it is the
same class of defect this audit reports in others' work: an instrument's limit
reported as a fact about the world.

\paragraph{The data.} The archive the link resolves to ($111{,}142{,}305$
bytes; its SHA-256 is in the repository) was downloaded by the first author,
read from the zip one line at a time, never extracted and never redistributed;
nothing from it enters the repository but sorted per-cell counts and tallies.
It holds 340 recordings, not 342, one per session, each in the toolkit's own
recording format: one frame record per action, the first being the
construction reset, and a final line carrying the toolkit's scorecard for the
session with per-play cumulative action counts at each level boundary and a
reset count. Every frame names the environment version we audited. The
parser would split a session at any full reset after its first action; none of
the 340 contains one, so recordings, sessions and plays all number 340. The
archive carries no participant identifier, so the report's ``first-time
players'' cannot be filtered for and every statistic below is over all 340.
Cells rest on between 2 and 40 plays (median 8), and 21 of the 183 on fewer
than 5.
Twenty-three recordings, all from March 2026, are from an earlier recorder
that wrote action names rather than ids and no construction line; 22 of them
complete at least one level.

\paragraph{They are plays of the shipped code.} We replayed every recording
action by action, with its recorded click coordinates, through the
environments in the repository, comparing rendered frame, state and level
after every step. All 340 reproduce state and level at every one of
$180{,}496$ steps: 321 under the shipped engine as configured, and 19 only
under the engine's \code{ONLY\_RESET\_LEVELS} switch
(Section~\ref{sec:humanengine}). 284 reproduce the rendered frame at every step under the
shipped engine as configured; 15 do so only under the switch; the remaining 41,
in four environments, agree on state and level throughout while $6{,}844$
frames differ in decoration --- for example, a pixel or two changing colour on most frames
of one environment, animation frames in another, a colour-swapped flash frame
in a third --- which the game logic does not read.

\paragraph{What the published baseline is.} For every play and every level it
completed we counted actions two ways: \emph{charged}, every action the scorer
counts, level resets included and the construction reset excluded; and
\emph{uncharged}, level resets excluded. Our charged count agrees with the
toolkit's own tally in the recording's final line on 316 plays; on the 22
old-recorder plays that complete a level the stored tally is exactly one below
ours at every level boundary, and the same recordings replay through the
shipped code with no shift, so it is the old tally that is off; two
\code{lp85} cards record a total and a level count but no boundaries. Over the
183 published integers, the upper median of the charged counts reproduces 165
exactly and four more within one; the upper median of the stored tallies
reproduces 166 --- it gains the two old-tally cells below and loses
\code{lp85} level 6, where the two boundary-less cards drop out and move the
median by one; the uncharged reading reproduces 101, with 42 within one. Lower
median, plain median and mean do far worse. The documented rule --- the upper
median of players' counts --- is the rule, and it charges resets. Because it reproduces over all 340 sessions,
it suggests either that every released session was a first play or that the
published baselines were not restricted to first-time players as the
documentation states; the release cannot distinguish the two. On the 72
cells where a human reset and the two readings differ, the published value
equals the charged reading on 65, the uncharged on one, and neither on six;
all seven cells that are not the charged reading lie in the three
environments the release does not reproduce.

\paragraph{Two level-1 cells carry the old recorder's tally.} The published
level-1 baselines for \code{cn04} and \code{tr87}, 29 and 54, equal the upper
median of the stored tallies and sit one below the upper median of the actions
those humans took by the current scorer's rule, 30 and 55. The discrepancy is
one action in two of the 183 cells; we record it because it is exact.

\paragraph{Three environments are not reproduced from the release.} Under the
rule that fits everywhere else, \code{g50t} reproduces four of seven cells,
\code{lp85} one of eight and \code{vc33} one of seven. Eight published values
occur in no released play of their cell at all. One is \code{cn04} level 1,
whose value is the old recorder's tally above. The other seven are two of
\code{g50t}'s --- level 4's baseline of 230 sits over two released plays that
took 31 and 52 actions --- and five of \code{vc33}'s seven. No statistic over
these recordings produces those seven numbers, so for those environments the
release is not the data the baselines were computed from. Splitting the
sessions by testing month does not help. One explanation is ruled out: the announcement's 342 against the
archive's 340 is a shortfall of two sessions, too few to account for seven
cells. Another we cannot test: a
revision of the environment after its human testing, which a version string
that matches would not exclude. We report this as not reproduced, not as
wrong.

\paragraph{Every baseline is at or above its optimum.} A human count bounds
the optimum from above. In all 183 cells the smallest human count is at or
below the published baseline, so every baseline is at or above its level's
optimum, the twelve the search could not resolve included. And in every cell
our search had bounded from below, no human beat the bound: the two
instruments agree. For the 165 reproduced cells this is close to automatic,
since the published value is itself a human count on the same engine; it
carries information for the other 18, whose values the release does not
reproduce but which still lie at or above a human count for their level.

\paragraph{How stable the baselines are.} Most cells rest on few plays, so we
asked how far each baseline would move under a different draw of players, a
question we registered before computing it. Resampling each cell's plays 2,000
times and recomputing the upper median each time, the median cell's 95\%
interval is 84.6\% as wide as its published value, and the published value
lies inside its own interval in 179 of the 183 cells. For an agent that
finishes a level in exactly its published baseline, and so scores 100 on it, the
median cell's interval corresponds to a level score between 50.6 and the cap of
115. The squared ratio passes the baseline's noise to the score with its
relative size doubled. For the 21 cells with fewer than 5 plays the resampling
understates the uncertainty, since it can only redraw the few counts it has. A
baseline measured from a median of eight people is a noisy yardstick, and the
per-level scores built on it inherit that noise.

\subsection{The engine the humans played}
\label{sec:humanengine}

The engine's \code{handle\_reset} performs a \emph{full} reset --- every level
re-cloned, score and level index zero, a new play opened on the scorecard ---
whenever its per-level action counter is zero and the game is not won, and a
level reset otherwise. That counter is zeroed by \code{set\_level}, which runs
at the start of every level and inside every level reset. So a \code{RESET}
issued as the first action of a level, or a second \code{RESET} straight after
a first, is a full reset. The engine's own switch \code{ONLY\_RESET\_LEVELS}
removes both cases: under it a reset never leaves the level except after a win.

The recordings show which reset semantics the humans played under. Across
the 340 plays, 29 resets fall at a zeroed counter, in 24 of the plays: 28 are
a reset straight after a reset --- a game over, a reset, then another --- and
one is a reset at the start of a level. On level 1 a full reset and a level
reset produce the same frame, state and level, so the replay can tell the two
settings apart only beyond level 1, where 24 of those 29 resets fall, in 19
plays. Those 19 recordings are exactly the ones that reproduce the humans'
state and level only under the switch; 15 of them also match every frame under
it, and four, in environments whose frames carry decoration, match on state
and level. Under the engine the humans played, each of the 29 resets cost one
action. Under the shipped default each of the 24 beyond level 1 would have been
a trip back to level 1. Whether the humans'
server ran the switch or the guard described below, the trace is the same;
what the replays establish is that their resets never left the level.

What that costs under the real harness we measured with the same
scripted-agent driver as Section~\ref{sec:forced}, on \code{ls20}: the 13-action
level-1 witness, one \code{RESET} as the first action of level 2, then the
45-action level-2 witness. Under the humans' setting the reset is a no-op
level reset; the play completes two levels and the environment scores
$10.714286$. Under the default the reset opens a second play at level 1, the
level-2 witness is spent on level 1, the card reads one level completed then
none, and the environment scores $3.571429$; the harness's own tracker still
records one level completed while the frame reports none, so its counter,
never reset and now at 46, is judged against level 1's budget. An agent that notices and re-clears level 1 recovers
the score at the price of a second play and 13 more charged actions.

Which of the two cases an agent can meet through the harness is settled by
the harness's own code. It offers \code{RESET} to the model on every turn
except the first and except immediately after a \code{RESET}, its own forced
one after a game over included, and it executes only actions it offered. So
the reset-after-reset case --- 28 of the 29 human events --- cannot occur
through the harness, and the forced reset of Section~\ref{sec:forced}, which leaves
the counter at zero, cannot be followed by a chosen one. A \code{RESET} as
the first action of a newly reached level is offered, and that is the case
the probe above measures.

Whether the server that scores agents runs the switch is not observable to
us. The toolkit's own API server carries a guard for exactly this case --- in
``competition mode'' a reset that would cause a full reset is not sent to the
game at all --- and the toolkit's changelog, for its own version 0.9.5 (the
toolkit package, distinct from the engine we pin at 0.9.3), fixes scorecards ``when
\code{ONLY\_RESET\_LEVELS} envvar was set'', so the Foundation knows the case
and handles it on at least two paths. The benchmarking harness sets neither,
and an agent run through it against the shipped environments meets the
default. The finding is the asymmetry that can be shown, on the shipped
engine: the reset semantics under which every baseline was gathered are not
the semantics the shipped default gives an agent run against those
environments. The budget finding of Section~\ref{sec:forced} is different in kind:
it is a property of the client, which by the code applies to online runs too.

\section{The scoring rule}
\label{sec:scoring}

\subsection{What the rule says}

Report v2 \S4.1 defines the level efficiency score for level $l$ of environment
$e$, with $a_{l,e}$ the agent's action count and $h_{l,e}$ the human baseline,
and gives the environment score as a level-index-weighted average subject to a
cap at the weighted fraction of levels completed. The human baseline is the
upper median (the report's term, defined in Section~\ref{sec:baselines}) of
first-time human players' action counts; report v1 said
second-best, and the change is documented in an April 2026 post
\cite{humandataset}.

\subsection{An inconsistency between the equation and the prose}

The report typesets the level score as
\[
S_{l,e} \;=\; \min\!\left(1.15,\ \frac{h_{l,e}}{a_{l,e}}\right)^{\!2},
\]
which caps the \emph{ratio} before squaring and therefore admits a per-level
score of $1.15^2 = 1.3225$. The same section's prose says the per-level
efficiency is capped ``at 1.15x human baseline'', which can be read as the
equation's ratio cap, and also that a level score lies between $0\%$ and
$115\%$, which cannot; the documentation says the score lies in that range; the
toolkit's changelog records the change from a $100\%$ cap to a $115\%$ cap; and
the shipped code applies \code{min(score, 115.0)} after squaring. Apart from the
equation, and one reading of that phrase, every artefact caps the score.

The difference is observable. Every planted trace has five levels with a
baseline of 10 on each, so the level weights sum to 15, and a level not listed
is left incomplete. Our planted trace \code{P2c} has one level solved
in $5$ actions against a baseline of $10$ and three solved in $12$: the shipped
scorer returns an environment score of \code{49.333333}, the prose reading
returns the same, and the equation reading returns \code{50.483333}. A level
solved in $9$ actions against a baseline of $10$ scores \code{115.0} in the
shipped code and $123.46$ under the equation.

\subsection{The action budget is not part of the scoring rule}

Report v2 \S4.3.1 states an action budget of five times the human baseline per
level, after which the agent is terminated. We find that budget implemented only
in the benchmarking harness, where every shipped model configuration sets a
multiplier of $5.0$ and the per-level budget is enforced in the agent's
termination test. The scorer applies no cutoff at all. Our planted trace
\code{P3b}, level 1 completed in $51$ actions against a baseline of $10$ and
the other four levels solved at baseline, is scored \code{93.589645} by the shipped scorer, where treating the budget as part
of the rule gives \code{0.000000}.

Neither the documentation's methodology page nor its competition-mode page
mentions the budget. The consequence is narrow but real: a scorecard computed by
the toolkit from a run that did not use the official harness --- an example
agent, a human session, a third-party agent --- is scored without the cutoff the
report describes.

We verified separately that the harness enforces the budget exactly. Across six
scripted runs against the toolkit in offline mode, driving the harness's own
control loop with a scripted model, no level ever exceeded its budget, the
harness's action counter and the scorecard's agreed on every run, and two runs
exited on the budget precisely at its boundary.

\subsection{The aggregation denominator}
\label{sec:denominator}

The scoring rule of Report v2 \S4.1 governs one environment. Above it the toolkit computes
an overall score, and there the denominator is the number of environments the
scorecard contains --- the number actually played --- rather than the number in
the set. The report defines the total as ``the sum of individual environment
scores divided by the total number of environments'', and the documentation's
methodology page describes it as the average of all game scores. Neither says
whether ``all'' means the environments in the set or the environments played, so
what follows is a discrepancy between two readings of an underspecified rule, of
which the implementation takes one; we report which reading the code takes and
what the other would give.

On a planted scorecard covering three environments of a set of 135, each
completed at exactly the human baseline on every level, the toolkit reports a
total of $100.0$ where the documented rule gives $2.222222222$, a factor of
$45.0$. Played over a whole set with every environment producing a card, the two agree exactly, so the difference is
precisely the partial-coverage case.

\begin{sloppypar}
Which runs this reaches is measured rather than assumed. Calling
\code{Arcade.get\_scorecard} under each operation mode, with a transport that
refuses to open a socket, identifies where the number comes from: \code{offline}
and \code{normal} compute it locally through this arithmetic, while
\code{online} and \code{competition} take a
\emph{remote fetch (server supplies the scorecard)}. The official benchmarking
harness constructs the toolkit with \code{operation\_mode} set to
\code{OperationMode.ONLINE}, so this arithmetic does not run for an official
run at all. The server-side scorer is not observable to us
and we make no claim about it, about the official leaderboard, or about any
published result.
\end{sloppypar}

The exposure is the local path, and it is not hypothetical. The toolkit's own
minimal example constructs \code{Arcade()} --- the default mode is
\code{normal} --- plays a single environment and prints \code{scorecard.score}.
Following the quickstart on one environment prints that environment's score
where the report's definition gives that score divided by the size of the set.

A second aspect sharpens it. Our first reading supposed that covering the whole
set is safe; on the local path it is not, because the divisor counts
environments that produced a card. Three environments played perfectly plus a
fourth that produced no card --- a failure, a skip, an id that did not match ---
gives a total of $100.0$ where the documented rule over four gives $75.0$. An
environment that produces a card and completes nothing, and one opened but never
played, both remain in the divisor and give $75.0$ correctly. The boundary is
between a card that exists and one that does not.

Relatedly, environment information is matched by full game identifier including
version. A scorecard recorded against one version and scored against a listing
carrying another scores that environment $0.0$, with the message
\emph{No Matching EnvironmentInfo found}. That moves a score down rather than
up, and a version refreshing between a run and its scoring would silently zero
it.

\subsection{Resets are charged to the agent, and to the humans}
\label{sec:resets}

A level reset increments both the reset count and the action count in the
scorecard, and the action falls on the level in progress, so it lowers that
level's efficiency. On an otherwise perfect five-level play, one reset during
each of levels 2 to 5 takes the environment score from $100$ to
$83.801652893$.

No document we found states whether a reset counts as an action. The methodology
page defines an action as a discrete interaction that affects the game state and
excludes internal operations, which does not settle it. The question had teeth
because human participants were explicitly permitted to reset a level at any
time, and the report observes that some ``reset levels after reaching a solution
in order to improve efficiency''. Had the published baselines been gathered
without charging human resets while agents are charged for theirs, every ratio
would be biased against the agent by an amount depending on how often each side
reset. The replays settle it (Section~\ref{sec:replays}): the published baselines are
the upper median of the toolkit's own per-play tallies, which charge every level
reset, and on the 72 cells (Section~\ref{sec:baselines}) where a
human reset and the two readings differ, the published value equals the
charged reading on 65, the uncharged reading on one, and neither on six. Humans paid for their resets as agents do. What remains is
a documentation gap: the rule is symmetric, and nothing says so.

\subsection{Resets the agent did not choose}
\label{sec:forced}

Section~\ref{sec:resets} concerns resets an agent chooses. The benchmarking harness
also issues them on its own: after a game over it returns \code{RESET} rather
than consulting the model, because a reset is the only way to continue. Those
resets are counted by the harness like any other action and, on the local path,
charged by the scorer to the level in progress.

Driving the harness's own control loop against the toolkit's fixture game (the
test environment the toolkit ships for its own unit tests), an
agent that loses a level three times and then solves it chooses 16 actions and
is counted 19, and the level is scored on all 19. Against that level's baseline
of 4, being scored on 19 rather than on the 16 the agent chose lowers that level's
score by 29\% relative to the score on 16 (the ratio $(16/19)^2$). That fall has a
counterpart in the baselines: a human who died also had to reset to continue,
and Section~\ref{sec:replays} shows those resets were charged too. What has no
counterpart is the budget, since the humans played without one.

A fall in a completed level's score is the mild consequence. The counter these
resets increment is also the one the harness stops on. The per-level budget is
$\lceil mb \rceil$ for the harness's multiplier $m$, which every shipped
configuration sets to $5.0$; the ceiling is the harness's own and is inert for
the integer baselines the benchmark publishes, so we write $\lceil 5b \rceil$
throughout. \code{is\_done} ends the level once the
counter reaches it; \code{choose\_action} returns a forced action before the
ordinary increment, but the routine that records the forced action increments
the same counter regardless. A run can therefore be cut off on an action the
agent did not choose.

We measured that as a paired counterfactual rather than arguing it from the
code. A scripted policy that dies once and chooses 8 actions runs twice against the same game with the same
budget, differing in one respect: in the second run the forced reset still
happens, still reaches the environment and is still counted by the scorecard,
but does not increment the budget counter. At a budget of 8 the shipped harness exits on the budget with the level
incomplete and the environment scoring \code{0.0}, while the counterfactual
completes the level and scores \code{1.316872428}. The agent's play is identical
in both. The reset it did not choose is the entire difference between a
completed level and a failed one.

The window is narrow, one budget value per death, but not accidental. Its width follows from the
mechanism as read from the code, and we wrote the prediction down before
measuring rather than assuming it. The counterfactual completes the level as soon as the budget covers the agent's
chosen actions, while the shipped harness additionally needs one action per game
over, so the set of budgets on which the two disagree should be exactly as wide
as the number of game overs. It is, in every case probed: no denial budgets for
an agent that never dies, \code{[8]} for one death, \code{[12, 13]} for two,
\code{[16, 17, 18]} for three. The general statement is simpler than the
demonstration. \emph{An agent's effective per-level budget is $\lceil 5b \rceil$
minus the number of times it died.} Each death quietly costs an action of
allowance on top of the progress it loses, and the exposure widens with every
one.

That demonstration uses the toolkit's fixture game and a budget we set, and a
fixture is not the benchmark. It reproduces on a real public environment at the
budget the harness derives from that level's published human baseline, which we
neither supply nor override. On \code{tu93} level 1, whose baseline of 19 gives
a budget of 95, an agent whose 95 chosen actions comprise one losing line and then a winning
one is cut off with the level incomplete and the environment scoring \code{0.0},
while the same policy with the forced reset uncharged to the budget counter
completes the level and scores \code{0.087046682}. The boundary appears on the same real level: at 93
and at 94 chosen actions both runs complete, and only at 95 does the reset
decide it. The losing line and the winning line are each replayed on the shipped
environment before use.

What the mechanism costs a policy that is not built to trip it is a separate
question, and we measured it rather than leaving the contrivance unqualified. A
seeded blind policy over each environment's advertised actions, clicks included
with coordinates inside the declared contract, was driven through the harness's
own loop twice per seed, once as shipped and once with the forced reset
uncharged; because the sequence is fixed in advance the two runs choose the same
actions as far as both go. Across 250 paired runs over all 25 environments and
ten seeds, the median share of charged actions that the agent did not choose is
$0.69\%$ and the maximum is $6.28\%$. The six environments that could be played
without clicks gave a maximum near 1\%; the largest observed tax is on a
click-only environment, so measuring only those six would have understated the
maximum. We report maxima, not a class comparison.

Where neither run completes a level, which is 234 of the 250, the extra actions
the uncharged run earns are exactly the forced resets the shipped run was
charged: the mechanism as an identity. In the remaining 16 it does not
hold, and should not, since those extra actions carry the uncharged run to a
different point in the game.

The outcome result is null, and it is as much a part of this finding as the
denial. No level's completion and no run's exit reason differ between the
charged and uncharged runs, in any of the 250. This is stronger than what six
environments could support: there the null could be explained by a policy that
never completes a level, whereas here 16 runs do complete one and the
outcome still never differs; we do not claim those 16 were near their
budgets, only that the trivial explanation no longer covers every run. The practical reach of this mechanism therefore
depends on an agent good enough to finish close to its budget.

We measured that reach on the play the baselines come from. Applying the
budget rule to every level a human completed in the released replays, with the
charged count against $\lceil 5b \rceil$ and the boundary as measured on
\code{tu93} level 1, the budget cuts off 29 of the 1,614 completions, and
none of them only because resets after a game over are counted: in every one
of the 29 the human's own chosen actions already exceed the budget. On
human-like play the budget binds rarely, and the counted resets never decide
it. What the mechanism can still do is end a run, since the harness stops the
whole environment when a level's budget is exhausted and forfeits the levels
after it; how often a capable agent finishes that close to its budget is not
measured here.

The sweep also cross-checks the exposure measurement across the whole set, and
the two agree in both directions: an environment where no loss inside the budget was
found never incurs a forced reset in any seed, and one where a loss was found
always does in some seed. Both instruments are random play, so this is a
consistency check rather than independent confirmation; the exhaustive result
on \code{ls20} level 1 below is the independent one. The comparison has to be made per environment across
its seeds rather than per run, because a single seed may complete the level
outright and never die; an earlier version of our own check compared one seed
and reported a disagreement that was not there.

The mechanism needs a game over, so we asked of all 183 levels of the 25 public
environments whether a level can be lost inside its own budget. Each shipped
environment is played from that level's start under the actions it advertises,
with click coordinates inside the declared contract, and the action sequence at
the first game over bounds the shortest losing line from above; the search is
capped at the budget, since a loss that takes longer is not exposure. Every
recorded line is replayed, coordinates included, to confirm it ends the game.

On 148 of the 183 levels, that bound falls below the budget, which proves the
level can be lost within it, and 23 of the 25 environments have at least one such
level. On the exposed levels the median recorded losing line is about a third of the
budget and the median number of deaths it affords is three; both are statistics
of upper bounds on the shortest loss. ``Deaths affordable'' is the budget divided
by the recorded line's length: how many copies of that line fit in the budget,
not a demonstration that the line repeats after each forced reset. On the remaining 35
no rollout died, and we report those as not established rather than as immunity,
since random play failing to lose says nothing about the shortest line.

The first version of this measurement sampled a fixed set of four actions for
every environment, including six that advertise only a click, and reported 17 of
25 first levels. Three of its recorded lines used actions the environment does
not advertise, two of them entirely; those actions are accepted and do advance
the state, but a line an agent is never told it may play is not a played loss.
Re-measuring under the advertised actions moves four environments and gives 19 of
25 at level 1, which is the figure we report. We record the correction because
the earlier number would otherwise have been published, and because it confirms
that an action outside \code{available\_actions} is accepted and counted: it
also advances the state and can end a level.

Where we can compare the two instruments they agree. On \code{ls20} level 1 an
exhaustive search puts the true shortest losing line at 129 against a budget of
110, so that level cannot be lost inside its budget at all, and the capped random
search correctly finds nothing there. That is why our first attempt to
demonstrate the denial there failed, and we record it rather than dropping it.

Applied to the real public set, one such reset costs a level between $0.035\%$ and
$3.33\%$ of its budget, and lowers an otherwise perfect completion (every level solved in exactly its
baseline, a level score of 100) by a
median, in level-score points, of \code{3.251814028} points across the 183 levels of the 25 public
environments. The worst case is the smallest baseline in the set, \code{sc25}
level 1 at $b = 6$, where one forced reset takes a perfect completion from
\code{100.0} to \code{73.469387755}.

Two things this is not. It is not a claim about how the server scores such a
reset, which is not observable to us. By the code, the budget itself is
enforced by the harness, the one its own README calls the official benchmarking
agent, and so would apply to online runs too; every run we made was offline. And it is not a design criticism: the harness has to
send something after a game over. What is absent is any statement to an entrant
that these resets happen, that they consume the per-level action budget, and
that they can end a level the agent would otherwise have completed. The humans' resets after a game over were forced by the game in the same way
and charged in the same way (Section~\ref{sec:resets}); the difference lies in
the budget alone.

Several negatives from the same instruments are worth recording, because each
rules out a way a client can distort its own score. A retried model call never
reaches the environment: the harness's retry path contains no call that does,
and across every scripted run the number of action-bearing environment calls
equals the number of actions counted; the construction reset below is not an
action. One intended action is sent exactly once: driven with a
transport that records every attempt and opens no socket, the remote wrapper
makes a single attempt under a connection error, a read timeout after the
request was sent, a server error and an empty response, so it never has the
server count two actions for one intent. A level that advances on the last
action the budget permits is not cut off. The frame cap subsamples evenly rather
than truncating and always keeps the last frame, so the settled frame is never
dropped; whether the humans who set the baselines saw more is not establishable
from the published material and we do not compare. And the reset the wrapper
issues when the environment is constructed goes on the wire but is counted by
neither the harness nor the local scorer, which treats it as beginning a play;
the server's treatment of it is unknown.

One further limit deserves mention without being a finding. Besides the action
budget, a run also stops on a wall-clock limit that the report does not describe.
It is uniform: none of the twelve shipped model configurations overrides it or
the frame cap, and all set the same action multiplier, so the base of twelve
hours for a single environment applies to every entrant. We entered the branch
by setting the limit to zero rather than by waiting, and confirmed it ends a run,
but at twelve hours per environment we could not show it binds.

\subsection{An unreached edge in the scorer}

The engine permits a game to advance the level counter twice within a single
action. Where that happens, the scorer pairs the $k$-th recorded level
transition with level $k$ and scores the final level as not completed even
though the play ends in a win: our five-level trace \code{P8}, every level solved at baseline, returns
\code{66.666667} with the first four levels credited and the fifth not,
where the documented rule gives \code{100.000000}.

We then asked whether any public environment can trigger it. Statically, all 25
sources contain exactly one \code{next\_level()} call site and none is inside a
loop. Dynamically, we found no such action anywhere: not in either complete
state graph of Section~\ref{sec:env}, not in the $1{,}245{,}427$ transitions
enumerated in Section~\ref{sec:sweep}, and not in the $479{,}040$ actions played
in Section~\ref{sec:play}. We report it as a robustness edge that a guard would
close, not as a defect with consequences.

\subsection{What the scorer gets right}

On 14 preregistered planted traces --- 13 of which have an oracle value, the
fourteenth exercising a case the rule leaves undefined --- the shipped scorer
agrees with our verified model of the prose rule on 12, both on environment
scores and on every per-level score, disagreeing on one: the double-advance edge
above. The
agreements include the $115$ cap, one-indexed level weights, the environment cap
at the weighted completed fraction, best-of-plays selection, and a level solved
in exactly the baseline. The model itself discharges 50 proof obligations with
no errors and no timeouts.

Two smaller documentation items: the code comments beside the cap still say
``max 100'', predating the change to $115$; and the published response schema for
\code{/api/games} omits \code{baseline\_actions}, a field the endpoint does
return and the toolkit reads.

\paragraph{A wrong prediction, disclosed.} We preregistered that the harness
would issue and count an initial reset, making it one action stricter than the
scorer on the first level. It does not: both toolkit wrappers issue that reset
inside environment construction, and neither the harness nor the scorecard
counts it. We added a further probe to exercise the boundary the original one
had missed, and report the corrected picture: the two counts agree.

\section{Environment models}
\label{sec:env}

\subsection{Method}

We selected one public environment by a rule fixed in advance: the fewest lines
of transition logic among environments with at least six levels and a mechanic
described in the report or documentation. Only \code{ls20} satisfies the last
clause --- the report shows its first level as a state graph and states a win
probability for it --- so it was selected even though eleven environments have
smaller sources. In effect the rule selects the one environment whose mechanic
the report describes; we state it as a rule because it was written before we
knew that. All 25 environment sources are identifier-obfuscated, so the
models are written from behaviour and structure rather than from names.

For each level we extract the constants from the shipped level object by script
(walls, start, goal, modifier tiles, energy pickups, step budget, decrement,
lives), hand-write the transition rules from the shipped step function at the
rule level, and generate a Dafny model. We then enumerate the level's reachable
state graph directly from the shipped game, and compare the compiled model
against the implementation on every transition of that graph and on 30 recorded
traces.

\subsection{Results}

Table~\ref{tab:ls20} summarises. States are counted under the enumerator's raw
key (Section~\ref{sec:state}). The three life-copies of that section contain every
state except the game-over states, 143 on level~1 and $1{,}291$ on level~2;
each level's single win state, the other absorbing state, lies inside them. Both models verify with no errors and no timeouts, and agree with the
shipped implementation everywhere. Because the environments' rules are stated
nowhere, the models are our rule-level reading of the shipped step function,
and agreement on every transition establishes that the implementation has no
behaviour outside the modelled rules on the reachable graph, not that those
rules are the intended ones.

\begin{table}[t]
\centering
\begin{tabular}{lrr}
\toprule
& \code{ls20} level 1 & \code{ls20} level 2 \\
\midrule
States in the complete reachable graph & 14{,}337 & 34{,}739 \\
Transitions & 56{,}772 & 133{,}788 \\
Minimum actions to complete & 13 & 45 \\
Published human baseline & 22 & 123 \\
Proof obligations discharged (0 errors) & 17 & 26 \\
Disagreements, model vs implementation, all transitions & 0 & 0 \\
Trace steps compared (30 traces) & 3{,}208 & 1{,}903 \\
Disagreements on traces & 0 & 0 \\
\code{RESET} returns to the level start & 500/500 & 500/500 \\
Actions advancing the level counter by two & 0 & 0 \\
\bottomrule
\end{tabular}
\caption{Two levels of the public environment \code{ls20}, audited against
machine-verified models of their transition rules.}
\label{tab:ls20}
\end{table}

The verified properties are: the level is winnable from its own start state,
witnessed by the shortest action sequence the enumeration found and replayed
through the model with every intermediate state asserted equal to the state the
shipped game is actually in; no action leaves the legal state space; and reset
restores the start. The first is a stronger statement than ``the run ends in a
win'': if model and implementation parted company anywhere along the winning
path, it would not verify.

\subsection{The report's stated win probability}

Figure~3 of the report shows the first level of \code{ls20} as a graph, notes
``three repeating states -- an artifact of the three-life mechanic'', and states
that the win probability for that level is ``exactly 1 in 355''. The term is not
defined. Our exhaustive enumeration gives a graph that is indeed three copies of
one structure, one per remaining life, and under the most natural reading
--- the probability that a uniformly random policy over the four movement
actions reaches the win before a game over, starting from the level's reset
state --- we compute \code{0.002813847150161035}, that is 1 in
\code{355.38533070027296}. That agrees to three significant figures; ``exactly''
is a rounding. Two other readings we computed, the reciprocal of the state count and the
winning fraction of states, both give 1 in $14{,}337$. We report this as a
definitional gap rather than an error: the environment agrees with the claim
under its natural reading.

\subsection{What counts as a state}
\label{sec:state}

A state count depends on the identity function used, and the choice is not
neutral. At the rule level the three life-copies are the same size up to a single
state on both levels: distinct states by remaining life are $4{,}732/4{,}731/4{,}731$ on level~1 and
$7{,}400/7{,}399/7{,}399$ on level~2. Our enumerator's raw key, which is deliberately
finer than the rendered frame, reports $7{,}400/13{,}024/13{,}024$ on level~2. The cause
is that a life loss re-appends previously consumed pickups to the end of the
level's sprite list, so a run that consumed one and a run that did not are
distinguishable by internal ordering while rendering identically; level~1 has no
pickups and is symmetric under both measures. This changes no transition and no
absorption probability, but it is a caveat on any published state count for
these environments, including the report's figure and our own: a frame hash, an
engine-state hash and a rule-level abstraction give three different answers.

\section{Which environments our enumerator can exhaust}
\label{sec:enumerable}
\label{sec:sweep}

Before enumerating anything we measured what is enumerable. Nineteen of the 25
public environments advertise the click action, whose $x$ and $y$ each range
over $0$ to $63$, giving at least $4{,}096$ distinct actions per state, most of them no-ops.
Exhaustive
enumeration is out of reach for those within any reasonable budget. The
remaining six --- \code{ls20}, \code{tr87}, \code{tu93}, \code{g50t},
\code{re86}, \code{wa30} --- advertise four or five simple actions.

This is a measured property of the advertised action space, taken from the
shipped games. It bounds enumeration over that space, which is what we do; a
method that quotients actions by their effect, or builds a partial or abstract
graph, is not bounded by it, and published agents do build graphs over these
environments.

For those six we enumerated level 1 by depth-first search, which holds only the
objects along one path and so is bounded by depth rather than by the width of a
search layer. Two completed within budget, \code{ls20} at $14{,}337$ states and
\code{tu93} at $2{,}609$; the other four stopped at a deterministic cap of
$150{,}000$ states, which establishes nothing about reachability either way and
is reported as such. Across $616{,}946$ states and $1{,}245{,}427$ transitions:
no level was shown unwinnable, \code{RESET} returned to the level's start on all
$3{,}000$ probes, and no action advanced the level counter by two.

\section{Reset behaviour across all 25 environments}
\label{sec:play}

Two of the three questions the enumeration answers need no state graph, so they
can be asked of the 19 environments enumeration cannot reach: does
\code{RESET} restore the state the level began in, and can one action advance
the level counter by two. We answer both by playing each environment at random
inside its own advertised action set, sampling click coordinates uniformly
within the declared bounds, and probing as we go.

We measure the reset question twice, because two different relations are at
stake. \emph{Frame identity} asks whether the rendered grid an agent observes
returns to the level's opening frame. \emph{Engine-state identity} asks whether
the game object does. Conflating them turns a property of the instrument into a
claim about an environment.

Across $18{,}663$ probes in all 25 environments, the rendered frame after a level
reset was identical to the frame the level began with every single time. Engine
state matched on $15{,}440$ of the same probes; in seven environments
(\code{cd82}, \code{ft09}, \code{m0r0}, \code{r11l}, \code{sb26}, \code{sc25},
\code{sp80}) some of the game object's own attributes, which account for every
probe on which engine state differed, survive a level
reset. Several are read back by the games' own code, so they could in principle
influence later behaviour, but none is visible in the frame at reset time.

We do not report this as a defect. ARC-AGI-3 states no rules for any environment
by design, so an environment that leaves an internal counter alone contradicts
nothing. It does bear on a documented assumption: the report's own graph
construction begins ``from the reset state of a selected level'', and published
graph-based agents identify states by hashing the frame. Frame identity and
engine-state identity are not the same relation, and in seven of 25 public
environments they disagree about whether a reset returns to a canonical state.
Anyone building on ``the reset state of a level'' should say which they mean.

The probe issues a reset only while the engine's per-level action counter is
nonzero, so none of its resets is the full reset of Section~\ref{sec:humanengine};
it reached level 2 in seven environments and no further.

Over $479{,}040$ actions in all 25 environments, no action advanced the level
counter by two and none moved it backwards outside a full reset.

\paragraph{Why these negatives are readable.} A detector that cannot fire proves
nothing. Both are tested against deliberately broken games: with state injected
that a reset does not restore, the probe reports a mismatch while the frame
still matches; with a reset made to leave a visible change, the count of frame
matches falls; and with the engine made to advance two levels in one action, the
detector catches it.

\section{Limitations}

A reader should hold the following in mind.

\begin{enumerate}
\item Everything here concerns the public toolkit, harness and documents at the
      pinned commits. The server-side scorer behind the official leaderboard is
      not observable to us and is not claimed to match. This bears hardest on
      Section~\ref{sec:denominator}: we report how a public artefact resolves an
      underspecified rule, not a
      claim about anyone's published result, and we do not compute what any
      particular reported score would have been.
\item ``No disagreement'' is over the fixtures and transitions we ran, chosen by
      us. It is not a proof of correctness.
\item Our Dafny models encode our reading of the report and the shipped source.
      A disagreement could always be the model being wrong, and we treated that
      as the default explanation throughout.
\item Two levels of one environment are audited with verified models. The
      state-graph results cover level 1 of six environments, two exhaustively;
      we could not enumerate the 19 click-based environments. Whether a
      level can be lost inside its budget is measured for all 183 levels, but
      that instrument is one-sided: it proves exposure where it finds a loss
      and proves nothing where it does not.
\item The random play probe won no game and mostly stayed on level 1, so its
      negatives cover the states it visited. The blind policy of
      Section~\ref{sec:forced} is likewise random: how often a capable agent
      finishes within a few actions of its budget, and so how far the budget
      finding reaches beyond random play, is not measured here.
\item The replay analysis reads the released recordings as they are. Three
      environments' baselines are not reproduced from them and are reported
      as such; the ratings file beside the archive was not used; and what the
      scoring server does with a reset at a zeroed counter is not observable.
\item A capped enumeration is a sample of a larger graph. Four of the six
      enumerable environments are represented by the first $150{,}000$ states
      a depth-first search happened to visit.
\item Several of our own errors were found during the work and are recorded in
      the repository's decisions file rather than quietly fixed: a state key
      that ignored non-\code{ls20} game state, two hashing errors that merged or
      spuriously split states, an off-by-one in the lower bound, a
      peak-memory guard that let one heavy environment mark later ones as
      capped, an exposure measurement that sampled actions six environments
      never advertise, a cross-check that compared one seed where it needed
      all of them, a stated HTTP status (429) that no artefact supported, and a
      shortener's refusal reported as the replays not being locatable. Each
      invalidated something that was then re-run or cut.
\item What a reader must still trust: Dafny and its JavaScript backend, Node,
      the vendored toolkit at the pinned commit, and our reading of the English
      in the report and documentation.
\end{enumerate}

\section{Related work}

This paper belongs to a line of work that audits a benchmark's reference data
rather than competing on it. Northcutt, Athalye and Mueller
\cite{labelerrors} find label errors in the test sets of ten widely used
datasets, frequent enough to make model selection by test accuracy unreliable.
Recht et al. \cite{imagenetv2} rebuild the CIFAR-10 and ImageNet test sets by
their original protocols and measure how far accuracy falls on the new ones.
Singh et al. \cite{leaderboard} show that undisclosed private testing distorts
the rankings of a public leaderboard. Raji et al. \cite{everything} question
what a benchmark can claim to measure, and Reuel et al. \cite{betterbench}
assess 25 AI benchmarks against 40 best practices and find large differences in
quality. Closest to our subject, Wei et al. \cite{humanbaselines} review 115
human baselines in foundation-model evaluations and find them rarely rigorous or
well documented, and in the Arcade Learning Environment, whose human-normalised
scores are the nearest precedent for ARC-AGI-3's, Machado et al.
\cite{machado} show how much results depend on evaluation protocol and Toromanoff,
Wirbel and Moutarde \cite{toromanoff} argue that claims of superhuman play
depend on which human baseline is used. Documentation standards such as datasheets \cite{datasheets} ask that
a dataset record how it was collected; the collection condition recovered in
Section~\ref{sec:humanengine}, the reset semantics the humans played under, is
the kind of fact such documentation exists to record. Our checks differ in kind: differential comparison
against the shipped implementation, exhaustive enumeration where the state space
allows it, and scripts whose outputs are registered against the paper's claims.
The machine-verified models cover the scoring rule and two levels of one
environment; the rest rests on testing, and we say which is which.


Rudakov, Shock and Cowley \cite{graphexplore} build hash-identified state graphs
over ARC-AGI-3 environments in order to \emph{solve} them, reporting no state
counts, no win probabilities and no verification. Rodionov \cite{worldmodels}
builds executable Python world models for all 25 public games, checked against
recorded observations rather than machine-verified, and does not audit
environments for defects. To our knowledge this is the first machine-verified
model of an ARC-AGI-3 environment's rules, the first exhaustive state counts
published for specific levels under a stated identity function, and the first check of the published human
baselines against proved optima and against the released human replays.

\section{Reproducibility and disclosure}
\label{sec:repro}

The audit repository, released with this paper at
\url{https://github.com/ibrahimmian36/arc-agi-3-audit},
contains the preregistrations, the decisions file recording
every deviation and every error of ours, a dated ledger, the generated models,
the scripts, the artefacts, a test suite, and a claims file with a checker that
verifies each number in this paper against the log that produced it. Environment
sources are not redistributed; they are fetched from the ARC API by a script.
The human replay archive is not redistributed either; a script downloads it
from the storage location the announced link resolves to, recorded in the
repository, and the repository carries its size, its hash,
and the per-cell counts derived from it, with no recording name, participant
identifier or frame. The three audited repositories were archived at Software
Heritage at intake.

This work was made public without prior private notice to the ARC Prize
Foundation. We record that plainly rather than implying a disclosure process
that did not happen. Two considerations decided it. Everything examined here is
already public, and every defect is in published material or in the client an
entrant runs on their own machine, so nothing in this paper tells a reader
something they could not have derived from the same public sources; there is no
exploitable secret to withhold. And the findings that change a number are
properties of the shipped client, which any entrant can verify against their own
run today. We have no reason to think the Foundation would disagree with any of
them, and the repository contains everything needed to check every one.


\section*{Broader Impact Statement}

This paper reports defects in the public artefacts of a benchmark that other
people's results are measured against, so the first effect of publishing it is
on a named third party's software and documentation. We have tried to make that
effect proportionate: every claim is bounded by what was measured, every number
is produced by a script that ships with the paper, the negatives are reported
beside the positives, and our own errors are recorded rather than removed. A
reader who wants to check any claim can run it.

The intended benefit is to the people who interpret ARC-AGI-3 scores. An
entrant who knows that the harness charges its own resets, and that the
baselines charge human resets too, reads a score more accurately than one who
does not; the benchmark's authors can correct what they choose to correct. The
method is not specific to this benchmark, and the same instruments would apply
to any interactive benchmark whose scoring code and reference data are public.

Two risks are worth naming. The first is that a finding of this kind can be
read as an attack on the people whose work it examines; we have written it as a
statement of what the shipped code does and does not do, and it is published
together with the evidence rather than as an assertion to be taken on trust. The
second is that describing how a budget is spent could be read as a way to score
higher. It offers no route to a score an agent has not earned: the charged
resets follow game overs, so the only way to avoid them is to play better, and
nothing here bears on the private environments, which we never touched. A third risk is
that a toolkit total over a subset of environments is reported as a benchmark
score, which Section~\ref{sec:denominator} shows would overstate it. The
human replay data is used without redistribution and reported only as
aggregates, so no participant is identifiable from this paper. No models were
trained and no new data was collected from people.

\section*{Use of AI tools}

Large language models (LLMs) were used throughout this work: to write and debug
the audit's code, and to run blind reviews of drafts, with models from two
different providers. We have used an LLM to generate portions of the text of
this paper, and we take full intellectual responsibility for all of it,
including every claim, number and reference. No claim rests on a model's
judgement: every number in the paper is produced by a script in the repository
and checked against the paper by the claims checker (Section~\ref{sec:method}),
every reference was checked at its source, and every disagreement the
instruments reported was read by a person before it was called anything.

\section*{Acknowledgements}

The ARC Prize Foundation publishes the toolkit, harness, agents, technical
report and documentation openly and under permissive licences. That is what
makes an audit like this possible from outside, and it is not the norm. We thank
them for it, and note that the great majority of what we checked was exactly
what it claimed to be.

\begin{thebibliography}{9}
\raggedright  % unbreakable identifiers otherwise stretch a justified line

\bibitem{arcagi3}
ARC Prize Foundation.
\emph{ARC-AGI-3: A New Challenge for Frontier Agentic Intelligence}.
arXiv:2603.24621, v2, April 2026.

\bibitem{humandataset}
G. Kamradt (ARC Prize Foundation).
\emph{Measuring Human Performance on ARC-AGI-3}.
\url{https://arcprize.org/blog/arc-agi-3-human-dataset}, 14 April 2026.

\bibitem{graphexplore}
E. Rudakov, J. Shock, and B. U. Cowley.
\emph{Graph-Based Exploration for ARC-AGI-3 Interactive Reasoning Tasks}.
arXiv:2512.24156, December 2025.

\bibitem{worldmodels}
S. Rodionov.
\emph{Executable World Models for ARC-AGI-3 in the Era of Coding Agents}.
arXiv:2605.05138, v2, June 2026.

\bibitem{dafny}
K. R. M. Leino.
\emph{Dafny: An Automatic Program Verifier for Functional Correctness}.
LPAR-16, 2010.

\bibitem{chollet}
F. Chollet.
\emph{On the Measure of Intelligence}.
arXiv:1911.01547, November 2019.

\bibitem{labelerrors}
C. G. Northcutt, A. Athalye, and J. Mueller.
\emph{Pervasive Label Errors in Test Sets Destabilize Machine Learning
Benchmarks}.
NeurIPS Datasets and Benchmarks Track, 2021.

\bibitem{imagenetv2}
B. Recht, R. Roelofs, L. Schmidt, and V. Shankar.
\emph{Do ImageNet Classifiers Generalize to ImageNet?}
ICML, PMLR 97:5389--5400, 2019.

\bibitem{leaderboard}
S. Singh, Y. Nan, A. Wang, D. D'Souza, S. Kapoor, A. \"Ust\"un, S. Koyejo,
Y. Deng, S. Longpre, N. A. Smith, B. Ermis, M. Fadaee, and S. Hooker.
\emph{The Leaderboard Illusion}.
arXiv:2504.20879, April 2025.

\bibitem{everything}
I. D. Raji, E. M. Bender, A. Paullada, E. Denton, and A. Hanna.
\emph{AI and the Everything in the Whole Wide World Benchmark}.
NeurIPS Datasets and Benchmarks Track, 2021.

\bibitem{betterbench}
A. Reuel, A. Hardy, C. Smith, M. Lamparth, M. Hardy, and M. J. Kochenderfer.
\emph{BetterBench: Assessing AI Benchmarks, Uncovering Issues, and
Establishing Best Practices}.
NeurIPS Datasets and Benchmarks Track, 2024.

\bibitem{datasheets}
T. Gebru, J. Morgenstern, B. Vecchione, J. W. Vaughan, H. Wallach,
H. Daum\'e III, and K. Crawford.
\emph{Datasheets for Datasets}.
Communications of the ACM 64(12), 2021.

\bibitem{humanbaselines}
K. Wei, P. Paskov, S. Dev, M. J. Byun, A. Reuel, X. Roberts-Gaal, R. Calcott,
E. Coxon, and C. Deshpande.
\emph{Position: Human Baselines in Model Evaluations Need Rigor and
Transparency (With Recommendations \& Reporting Checklist)}.
ICML, PMLR 267:82265--82325, 2025.

\bibitem{machado}
M. C. Machado, M. G. Bellemare, E. Talvitie, J. Veness, M. Hausknecht, and
M. Bowling.
\emph{Revisiting the Arcade Learning Environment: Evaluation Protocols and Open
Problems for General Agents}.
Journal of Artificial Intelligence Research 61:523--562, 2018.

\bibitem{toromanoff}
M. Toromanoff, E. Wirbel, and F. Moutarde.
\emph{Is Deep Reinforcement Learning Really Superhuman on Atari? Leveling the
Playing Field}.
arXiv:1908.04683, August 2019.

\end{thebibliography}

\appendix
\section{Protocols, identities and resources}
\label{app:protocol}

Everything in this appendix is also in the repository; it is repeated here so
that each result can be interpreted from the paper alone.

\subsection{Environment identity}

Every public environment is identified by the game id the ARC API returns,
which carries a version hash (\code{ls20-9607627b}), and by the SHA-256 of the
source file the API served on 4 September 2026, of which Table~\ref{tab:envs} gives the
first 16 hexadecimal digits. The sources are not redistributed; a reader
who fetches the same ids can confirm they hold the same files.

\begin{table}[h]
\centering\small
\begin{tabular}{lrl}
\toprule
Game id & Levels & Source SHA-256 (first 16) \\
\midrule
\code{ar25-0c556536} & 8 & \code{b965d2faaa4a4d91} \\
\code{bp35-0a0ad940} & 9 & \code{e9aecb52c629c3e7} \\
\code{cd82-fb555c5d} & 6 & \code{dafeabff60a3a99b} \\
\code{cn04-2fe56bfb} & 6 & \code{29977e145a0374b5} \\
\code{dc22-fdcac232} & 6 & \code{74abe3523bd6bc8a} \\
\code{ft09-0d8bbf25} & 6 & \code{aa5b54f48a29da9f} \\
\code{g50t-5849a774} & 7 & \code{6f21f2065744e484} \\
\code{ka59-38d34dbb} & 7 & \code{3e8341f4c317ceb9} \\
\code{lf52-271a04aa} & 10 & \code{a6b01973a0905a44} \\
\code{lp85-305b61c3} & 8 & \code{7d10baec7dd29d07} \\
\code{ls20-9607627b} & 7 & \code{298c810da2850d55} \\
\code{m0r0-492f87ba} & 6 & \code{fc3954236f712759} \\
\code{r11l-495a7899} & 6 & \code{95d58ebaa8c758e9} \\
\code{re86-8af5384d} & 8 & \code{cf3c1520e17f0b70} \\
\code{s5i5-18d95033} & 8 & \code{341f85bac1fb5ca7} \\
\code{sb26-7fbdac44} & 8 & \code{dbb4877853a8d30f} \\
\code{sc25-635fd71a} & 6 & \code{871a507cf340d1e4} \\
\code{sk48-d8078629} & 8 & \code{b6ca05eed7ba81fe} \\
\code{sp80-589a99af} & 6 & \code{2306d62311a3cc02} \\
\code{su15-1944f8ab} & 9 & \code{a5f91f7c963d6ca6} \\
\code{tn36-ef4dde99} & 7 & \code{43c30052cdeb2300} \\
\code{tr87-cd924810} & 6 & \code{3274657a6499af5c} \\
\code{tu93-0768757b} & 9 & \code{80e41888f9f7b1a0} \\
\code{vc33-5430563c} & 7 & \code{8afa9f55054b2a24} \\
\code{wa30-ee6fef47} & 9 & \code{0c8b63caee2d092d} \\
\bottomrule
\end{tabular}
\caption{The 25 public environments as audited.}
\label{tab:envs}
\end{table}

\subsection{State identity}
\label{app:key}

The enumerator's raw key is a 128-bit BLAKE2b digest of, in order: every sprite
of the current level in list order (name, tags, position, visibility, rotation
and pixel array; order matters because collision scans iterate in order and can
break early); the engine's state name, score and level index; and every
instance attribute of the game object in sorted name order, serialised by
content to depth six, with dictionary keys and set members ordered by their own
content digests rather than by memory address. Excluded are the five per-action
bookkeeping fields (\code{\_action}, \code{\_action\_complete},
\code{\_action\_count}, \code{\_full\_reset}, \code{\_next\_level}) and the two
level containers (the other levels, which no action mutates, and the immutable
templates from which a reset re-clones a level). The key therefore covers every
mutable field of the game object other than that bookkeeping, so equal keys are
identical engine states for every action the enumeration takes. The enumeration
issues no \code{RESET}; a reset's effect also depends on the excluded per-level
action counter (Section~\ref{sec:humanengine}).
The engine is deterministic given its state: 23 of the 25 sources use no
randomness, and the two that do (\code{lf52}, \code{tr87}) seed it from the
level index. Equal keys thus have equal futures under those actions, and an enumeration over keys is
exhaustive with respect to the engine's own state. The two hashing errors of
Limitation~8 were in the canonicalisation --- ordering by memory address, and
hashing numeric scalars by type rather than value --- and were caught by an
address-independence test before any run reported here. Engine-state identity
in Section~\ref{sec:play} is equality of this key; frame identity is equality
of the rendered grid. The rule-level counts of Section~\ref{sec:state} use a
hand-written abstraction of the level (position, orientation, colour, shape,
lives, steps, goals and consumed pickups) and are reported for comparison only.

\subsection{Exposure protocol (Section~\ref{sec:forced})}

For each of the 183 levels the shipped game is started at that level. Actions
are drawn uniformly from the environment's advertised set; a click carries
coordinates drawn uniformly from $0$ to $63$ in each axis. A rollout stops at
\code{GAME\_OVER}, at a win or level advance (not a losing line), or after
$\lceil 5b \rceil$ actions, since a loss that takes longer is not exposure.
Eight rollouts are run per level, each with its own generator:
Python's \code{random.Random} seeded with the string \code{"<game>:<level>:0:<i>"} for
$i=0,\dots,7$. The shortest recorded losing line is kept with its coordinates
and replayed on the shipped game to confirm it ends in \code{GAME\_OVER}. The
per-level outputs are the line, its length, how each rollout ended, and the
budget.

\subsection{Blind-policy sweep (Section~\ref{sec:forced})}

For each environment and each of ten seeds, a script of $\sum_l \lceil 5 b_l
\rceil + 8$ actions (8 spare, so the script outlives any run) is drawn in advance, uniformly over the advertised actions
with click coordinates as above, from \code{random.Random} seeded with
\code{"<game>:<seed>"}. The harness's own control loop runs it against the
shipped environment offline, with the budgets the harness derives from the
published baselines. The paired run uses the identical script; its only
difference is that the increment of the per-level budget counter for a forced
reset is undone immediately after the harness applies it, with the reset itself
still sent and still counted by the scorecard. A script of that length cannot
be exhausted, because chosen actions never exceed the sum of the budgets. The
per-run outputs are the chosen and forced counts, the scorecard's action and
reset counts, the levels completed, the exit reason and the environment score.

\subsection{Data used and released}
\label{app:data}

\paragraph{Used, not redistributed.} The 25 public environments, fetched from
the ARC API by a script in the repository and identified in
Table~\ref{tab:envs}; the toolkit, harness and example agents at the pinned
commits, all MIT-licensed; and the ARC Prize Foundation's human replay archive,
downloaded from the location the announcement's link resolves to
(Appendix~\ref{app:replay}). None is redistributed. Readers obtain each from its
source, and the repository records the identity of the copy we used (version
hash, commit or SHA-256) so that a reader can confirm they hold the same one.

\paragraph{Released.} The audit repository (Section~\ref{sec:repro}), under the
Apache 2.0 licence: the scripts, the generated models, the preregistrations,
the decisions file and the ledger, and every artefact from which the paper's
numbers are read, including the per-cell sorted action counts derived from the
replays. The derived counts carry no recording name, participant identifier,
timestamp or frame.

\paragraph{Licence of the human replays.} The archive carries no licence file;
the Foundation's announcement describes the dataset as open-sourced. We read it
and report only aggregates derived from it.

\paragraph{Intended uses.} The derived counts support checking the published
baselines and any statistic over them; the scripts support re-running every
measurement in this paper against the same public artefacts or later versions.

\paragraph{Hosting and maintenance.} The repository is public on GitHub, and the
audited repositories were archived at Software Heritage at intake. A test suite
checks every figure in the paper against the artefact it comes from. A number
is corrected by re-running the script that produced it, and errors are recorded
in the decisions file rather than silently fixed. The authors maintain the
repository.

\paragraph{Hardware.} Every measurement ran on one laptop (Apple M2 Pro, 16 GB of
memory, macOS); the caps and peak memory of each instrument are stated in this
appendix.

\subsection{Enumeration and search caps}
\label{sec:caps}

The two complete \code{ls20} graphs were enumerated breadth-first under caps of
$400{,}000$ states, $1{,}500$ seconds and $3{,}500$~MB of resident memory, and
completed well inside them. The level-1 sweep of Section~\ref{sec:enumerable} ran
depth-first with caps of $150{,}000$ states, $1{,}200$ seconds and $2{,}500$~MB;
the baseline search of Section~\ref{sec:baselines} breadth-first with $400{,}000$
states, $120$ seconds and $2{,}500$~MB per level; the reset probes of
Section~\ref{sec:play} with $20{,}000$ actions and $60$ seconds per environment at seed
$0$. Every artefact records the peak resident memory of the run that produced
it. All runs were made on one laptop.

\subsection{Replay protocol (Section~\ref{sec:replays})}
\label{app:replay}

Archive \code{arc\_agi\_3\_public\_demo\_human\_testing.zip}, $111{,}142{,}305$
bytes, SHA-256 beginning \code{99a32ffc}, read from the zip one line at a time.
Recording $=$ one JSON object per line; a line with \code{action\_input} is a
frame after one action; the final line is the session's scorecard. A play is
split at a frame flagged \code{full\_reset}. Level $k$, counting from zero, is
credited when \code{levels\_completed} first exceeds $k$; a level never
completed contributes no count. Charged count $=$ actions at that level, level resets included, the
construction reset (the first record, when it is a \code{RESET}) excluded;
uncharged $=$ the same without level resets. Upper median of $n$ sorted values
$=$ the value at index $\lfloor n/2 \rfloor$. A cell is decisive when at least
one human reset there and the charged and uncharged upper medians differ.
Replay: the environment is constructed as the toolkit constructs it, each
recorded action applied with its recorded \code{x}, \code{y}; a recording is
classed by the weakest engine setting under which every frame, state and level
matches, and ``decoration only'' requires state and level to match at every
step under the switch. A trap event is a \code{RESET} while the engine's
per-level counter is zero and the game is not won, split by whether the
previous action was also a \code{RESET}. Peak resident memory stayed under
300~MB in both the baseline script and the replay script.

\subsection{Tools and proof obligations}

Dafny 4.11.0 with its JavaScript backend, Node 24.16, Python 3.12.13,
\code{arcengine} 0.9.3, and the three audited repositories at the commits of
Section~\ref{sec:scope}. The supplementary archive records the commit it was built from in its
\code{COMMIT} file. The scoring model's 50 obligations comprise the definitions of
the prose and equation readings and lemmas S1--S6: both caps bound the level
score (S1), a level solved in exactly its baseline scores 100 (S2), the level
score is monotone non-increasing in the action count (S3), the environment score
is capped at the weighted completed share (S4), every level at baseline gives
100 (S5), and nothing completed gives 0 (S6), with the arithmetic lemmas they
rest on. Each environment model discharges E1 (the level is winnable by the
witness, with every intermediate state asserted equal to the shipped game's),
E2 (no action leaves the legal state space) and E3 (reset restores the start).

\end{document}
